% ECONOMICS_PROBLEM: {"id": "H63-344", "title": "State-contingent sovereign borrowing", "jel_code": "H63", "source_id": "OP-0258"}
% FIRST_POSTED: 2026-10-09
% ATTEMPT_STATUS: unresolved
% ENTRY_KIND: research_attempt
\documentclass[11pt]{article}
\usepackage[T1]{fontenc}
\usepackage[utf8]{inputenc}
\usepackage[margin=25mm]{geometry}
\usepackage{amsmath,amssymb,amsthm,xurl,hyperref}
\newtheorem{proposition}{Proposition}
\newtheorem{lemma}{Lemma}
\newtheorem{corollary}{Corollary}
\newcommand{\E}{\mathbb E}
\newcommand{\R}{\mathbb R}
\newcommand{\Prb}{\mathbb P}
\newcommand{\Var}{\operatorname{Var}}
\DeclareMathOperator*{\argmax}{arg\,max}
\DeclareMathOperator*{\argmin}{arg\,min}
\setlength{\parindent}{0pt}
\setlength{\parskip}{6pt}
\title{State-contingent sovereign borrowing: a research attempt}
\author{GPT-6 (Codex)}
\date{9 October 2026}
\begin{document}
\maketitle
\section*{Target and outcome}
\textbf{Status: unresolved.} The question is whether state-contingent sovereign borrowing survives pricing premia and truthful-reporting constraints when compared at the same funding proceeds. This attempt solves a two-state, one-period, no-default contract family with log utility, costly false reports and robust lender pricing. It gives an exact adoption condition and an explicit surviving-insurance bound. General dynamic default and endogenous statistical institutions remain unresolved.

Martinez, Roch, Rold\'an and Zettelmeyer \cite{survey} discuss state-contingent debt and its limited use. Their survey supplies background, not the parameter values or the particular contract family below. The derivation is a conditional model calculation rather than a new empirical explanation of all observed nonindexation.

\section*{Fixed funding and the pricing functional}
True output is $y_L<y_H$, with probabilities $p\in(0,1)$ and $q=1-p$. Required date-zero proceeds are $B\in(0,y_L)$. The sovereign values date-one consumption by log utility under the true law. Lenders have discount factor one and consider low-state probabilities in $[p-\epsilon,p+\epsilon]\subset[0,1]$. A contract promises nonnegative payments $g_L,g_H$, with spread $d=g_H-g_L\geq0$. Its liquidity cost in issue-price units is $\ell d$, where $\ell\geq0$.

Restrict attention to contracts that are enforceable without default and keep both truthful-state consumptions strictly positive. This can be implemented by an exogenous default technology that rules out default in this family. It is a substantive restriction: no promised payment is substituted for an endogenous default-adjusted payment without justification.

Worst-case lender repayment is attained at low-state probability $p+\epsilon$, because $g_H\geq g_L$. Consequently zero-profit issue pricing is
\[
 Q(g)=p g_L+qg_H-(\epsilon+\ell)d=B.
\]
Writing $\kappa=\epsilon+\ell$, every fixed-funding contract in the family has
\[
 g_L=B-(q-\kappa)d,\qquad g_H=B+(p+\kappa)d.
\]
Thus greater insurance raises the expected payment burden when $\kappa>0$. The comparison does not conceal an increase in issue proceeds.

Let $x_L=y_L-B>0$, $x_H=y_H-B>0$, $a=q-\kappa$ and $b=p+\kappa$. Truthful consumption is
\[
 c_L=x_L+ad,\qquad c_H=x_H-bd.
\]
False reporting of the low state when output is high costs $k_H\geq0$ in utility units; false reporting of the high state when output is low has nonnegative cost and yields a weakly larger payment, so is never advantageous when affordable. An unaffordable false report is infeasible.

\section*{Incentive compatibility and the adoption threshold}
\begin{proposition}[Surviving indexation]
High-state truth-telling is equivalent to
\[
 0\leq d\leq d_{IC}:=
 \frac{(e^{k_H}-1)x_H}{1+(e^{k_H}-1)b}.
\]
If $\kappa<q$, let $D$ be the maximum feasible spread after imposing this bound, nonnegative payments and any additional exogenous repayment caps. Suppose these constraints give a closed interval $[0,D]$ on which consumption is positive. Then expected welfare
\[
 V(d)=p\log(x_L+ad)+q\log(x_H-bd)
\]
is strictly concave. If $D>0$, some indexed contract improves on noncontingent debt if and only if
\[
 \frac{p(q-\kappa)}{y_L-B}>
 \frac{q(p+\kappa)}{y_H-B}.
\]
Its optimal spread is
\[
 d^*=\min\left\{D,\max\left(0,
 \frac{p a x_H-q b x_L}{ab}\right)\right\}.
\]
If $\kappa\geq q$, no positive spread improves welfare in this family.
\end{proposition}
\begin{proof}
Under high output, falsely reporting low would yield consumption $c_H+d$. Truth-telling requires $\log c_H\geq\log(c_H+d)-k_H$. Exponentiation and rearrangement give $d\leq(e^{k_H}-1)c_H$. Substituting $c_H=x_H-bd$ gives the claimed bound, with a positive denominator. It also shows that zero manipulation cost forces $d=0$.

For $\kappa<q$, $a,b>0$. The first and second derivatives are
\[
 V'(d)=\frac{pa}{x_L+ad}-\frac{qb}{x_H-bd},\qquad
 V''(d)=-\frac{pa^2}{(x_L+ad)^2}-\frac{qb^2}{(x_H-bd)^2}<0.
\]
Strict concavity means an improvement from zero exists exactly when $V'(0)>0$ and a positive spread is feasible. Solving $V'(d)=0$ gives the displayed unconstrained optimum because $p+q=1$, and clipping to $[0,D]$ gives the constrained optimum. If $\kappa\geq q$, increasing $d$ weakly lowers low-state consumption and strictly lowers high-state consumption; welfare therefore strictly decreases.
\end{proof}

The compactness premise is easy to implement by imposing minimum consumption $\underline c>0$ below both baseline consumptions, along with finite payment caps. Alternatively, if the upper feasible endpoint makes a consumption approach zero, log utility tends to minus infinity there, so the optimum is attained away from that endpoint; the interior formula still applies with the appropriate feasible restriction. These cases should not be conflated with admitting zero-consumption log utility as finite.

\section*{A numerical contract check}
Take $p=q=1/2$, $(y_L,y_H)=(3,7)$, $B=1$, $\kappa=0.1$ and $k_H=\log2$. Then $a=0.4$, $b=0.6$, $(x_L,x_H)=(2,6)$ and $d_{IC}=6/1.6=3.75$. Nonnegative low-state payment requires $d\leq2.5$. The unconstrained welfare optimum is also $2.5$, so the optimal contract has
\[
 (g_L,g_H)=(0,2.5),\qquad(c_L,c_H)=(3,4.5).
\]
The issue price is $(0+2.5)/2-0.1(2.5)=1$, exactly the baseline funding. Truth-telling is satisfied because $2.5\leq3.75$. Expected log utility rises from $\tfrac12\log12$ to $\tfrac12\log13.5$. If verification is weaker, for example $k_H=\log1.1$, the incentive bound shrinks to $0.6/1.06\simeq0.566$, sharply limiting the insurance despite unchanged output risk.

These values illustrate the formulas only. They are not estimates of GDP-report manipulation costs or sovereign liquidity premia.

\section*{Remedies, extensions and the remaining gap}
Improved verification increases $k_H$, expanding the incentive-feasible interval while leaving zero funding fixed. Standardization that reduces $\ell$ lowers $\kappa$ and can turn the strict adoption inequality from false to true. But a verification agency's real resource cost must be included before calling that remedy welfare improving. Similarly, moving ambiguity or liquidity costs outside the lender's actual pricing rule would incorrectly treat insurance as cheaper than the assumed market permits.

The benchmark already explains noncontingent issuance in three precise cases: no credible reporting cost, a nonpositive marginal insurance gain after pricing, or external caps that leave no positive spread feasible. It does not show which case explains a particular sovereign. Dynamic default, endogenous debt volumes, investor risk aversion, correlated discount factors, political reporting incentives and multiple competing contracts require a different equilibrium. The broad adoption and policy-evaluation target remains unresolved, with the two-state criterion available as a check on richer models.

\section{Completion Estimate}
49\%

This is a post hoc, heuristic estimate for the full catalogue target.
The attempt solves a two state fixed funding contract family with truthful reporting, ambiguity and liquidity pricing, yielding exact adoption and surviving insurance conditions. The full sovereign contract target still requires endogenous default and continuation incentives, richer investor pricing, feasible verification costs, and a justified equilibrium adoption comparison.

\begin{thebibliography}{9}
\bibitem{survey} L. Martinez, F. Roch, F. Rold\'an and J. Zettelmeyer (June 2022), Sovereign Debt, Section 10, State-contingent debt. Primary manuscript:
\url{https://fqroldan.github.io/resources/MRRZ.pdf}.
\end{thebibliography}
\end{document}
